\chapter[Extracellular fields and inverse problems]{From membrane currents to extracellular electrodes}
\label{ch:electrode-physics}
An extracellular electrode does not measure a neuron's voltage
in isolation. It measures a difference of electric potential
created by spatially distributed ionic currents, tissue geometry,
conductivity, referencing, and acquisition electronics.

\section{Electrochemical driving forces}
For an ideal dilute permeant ion with valence $z\ne0$, equality
of electrochemical potentials on the two sides of a membrane gives
\[
E_{\mathrm{ion}}=\frac{RT}{zF}
  \log\frac{c_{\mathrm{out}}}{c_{\mathrm{in}}}.
\]
This is the Nernst equilibrium potential under its concentration
and temperature assumptions. The Goldman--Hodgkin--Katz formula
extends the constant-field model to multiple monovalent ionic
species with permeability parameters; it is not equivalent to
measuring those permeabilities from a single LFP.
In the Hodgkin--Huxley description the sodium current is modeled
as $I_{\mathrm{Na}}=\bar g_{\mathrm{Na}}m^3h(V-E_{\mathrm{Na}})$;
gating variables obey voltage-dependent kinetic equations.
The passive cable approximation has the schematic form
\[
c_m\partial_t V(x,t)=g_a\partial_{xx}V(x,t)
   -g_m(V-E_L)+i_{\mathrm{ext}}(x,t),
\]
with coefficients per unit length and consistent units.
Cable and membrane models determine transmembrane sources; they
are not by themselves electrode forward models.

\section{Quasistatic field equation}
At electrophysiological frequencies and spatial scales where
inductive effects are negligible, write $\mathbf E=-\nabla\phi$.
Let impressed current $\mathbf J_p$ and Ohmic return current
$-\boldsymbol\sigma\nabla\phi$ sum to a divergence-free total
current. Current conservation implies
\[
\nabla\!\cdot(\boldsymbol\sigma\nabla\phi)
       =\nabla\!\cdot\mathbf J_p,
\quad\text{or}\quad
-\nabla\!\cdot(\boldsymbol\sigma\nabla\phi)=q,
\quad q=-\nabla\!\cdot\mathbf J_p.
\]
The tensor $\boldsymbol\sigma$ may be anisotropic and spatially
inhomogeneous. To solve the equation, the domain, conductivity,
boundary conditions and electrode reference must be specified.
For an infinite isotropic medium, a point source of strength $I$
located at $\mathbf r_0$ gives, away from that point,
\[
\phi(\mathbf r)=\frac{I}{4\pi\sigma
  |\mathbf r-\mathbf r_0|}.
\]
Superposition yields a discrete source model $\mathbf v=G\mathbf s$
when the forward physics is linear, with lead-field entries obtained
by integrating the Green kernel over chosen source basis functions.
A dipole arises as the leading far-field contribution when net
source current cancels; line-source models integrate along a dendrite.

\begin{proposition}[Uniqueness for fixed Dirichlet conditions]
Let $\Omega$ be a bounded connected region and
$\boldsymbol\sigma$ uniformly positive definite. If two weak
solutions of the same elliptic equation have equal Dirichlet
boundary values, they agree (almost everywhere).
\end{proposition}
\begin{proof}
Their difference $w$ satisfies the homogeneous equation and
zero boundary values. Multiply by $w$, integrate by parts,
and use the boundary condition:
$\int_\Omega\nabla w\cdot\boldsymbol\sigma\nabla w\,dx=0$.
Positivity implies $\nabla w=0$; the boundary forces $w=0$.
For pure Neumann conditions, uniqueness holds only up to an
additive constant, with a compatibility condition on source flux.
\end{proof}

\section{Why the inverse need not be unique}
Suppose an electrode array measures $\mathbf v=G\mathbf s$.
If $G\mathbf h=0$ for nonzero $\mathbf h$, then
$G(\mathbf s+\mathbf h)=G\mathbf s$. This is exact
finite-dimensional nonidentifiability. In the continuum, smoothing
operators can also be ill-conditioned or compact, so tiny
measurement errors may cause enormous source errors even when
uniqueness holds in a restricted model.
\emph{Hadamard well-posedness} requires existence, uniqueness
and continuous dependence of solutions on data. A linear operator
is \emph{compact} if it maps bounded sets into sets with compact
closure; its singular values in infinite-dimensional inverse problems
can converge to zero.
For singular system $Gv_j=\sigma_j u_j$ and data $b$,
the formal inverse contains $\langle b,u_j\rangle/\sigma_j$
and magnifies noise when
$\sigma_j\to0$. A classical Picard-type condition controls the
squared normalized coefficients for compatible data.
Tikhonov regularization instead solves
\[
\widehat{\mathbf s}_\lambda
=\arg\min_{\mathbf s}
\{\|G\mathbf s-\mathbf v\|_2^2
 +\lambda\|\mathbf s\|_2^2\}
=(G^\top G+\lambda I)^{-1}G^\top\mathbf v.
\]
\begin{example}[MATLAB and interpretation]\sloppy
\path{field_potential_forward.m} receives $G$;
it does \emph{not} compute tissue conductivity or the Green kernel.
\path{field_potential_inverse_tikhonov.m}
returns one regularized source estimate, not a unique membrane-current
pattern. \path{current_source_density.m} needs
spatial geometry and differentiated potentials; references and
volume conduction remain crucial.
\end{example}